\subsection{Fonction structurelle de \texorpdfstring{\(K\) et \(\alpha\)}{K et alpha} dans l'incidence exceptionnelle}

Le lien entre la détermination numérique et la réalisation exceptionnelle
admet une description opératoire. Dans la lecture arithmétique, \(K\)
agit dans la normalisation des registres régionaux. Dans la lecture
d'incidence, ses composantes sélectionnent le support
\[
 B_K=\{j:K_j\geq729\}=\{4,6,9,10\}.
\]
Le même registre ordonné participe donc à deux fonctions
spécifiées : une composition positionnelle et une sélection d'incidence.
Sa réalisation rationnelle conserve le registre lorsque la base,
la longueur et l'origine sont fixées, tandis que le passage au support retient
l'information nécessaire au drapeau combinatoire.

Le registre signé antérieur à la normalisation de \(\alpha\) fournit
l'hexade \(H_\alpha^-\). Les intersections de cette hexade avec les
supports régionaux et avec l'incidence d'APP déterminent le point
marqué. Ce point individualise le vecteur de norme \(54\) utilisé
dans le voisin du réseau. Ainsi, l'orientation qui intervient dans la
normalisation arithmétique agit à nouveau dans la sélection géométrique.
La relation ne requiert pas de déduire un réseau à partir d'une approximation
décimale : elle utilise les données signées antérieures à son évaluation.

La réalisation exceptionnelle conserve successivement différents types de
relation. Le code conserve les combinaisons linéaires ; le plan conserve
les incidences entre supports ; le recollement conserve les classes discriminantes ;
le réseau conserve sa forme bilinéaire ; l'algèbre de vertex conserve la
graduation et les produits correspondants. Leurs groupes d'automorphismes
se définissent par rapport à ces structures. L'apparition du Monstre à
la fin de la construction exprime une propriété de l'algèbre
\(V^\natural\) obtenue par la procédure citée, et Moonshine
relie les traces graduées de cette action à des fonctions modulaires.

Cette hiérarchie fournit une signification mathématique concrète à la
fonction de liaison attribuée à \(K\) et \(\alpha\). Le premier est un
registre entier réversible avec des réalisations positionnelles et
d'incidence ; la seconde est la coordonnée de compatibilité arithmétique
dont le registre orienté intervient également dans l'incidence marquée.
Les définitions acquièrent leur contenu par ces opérations et par les
identités exposées dans ce chapitre. Le développement de l'incidence
est une partie substantielle de l'argument de l'article : il montre quelles relations
subsistent lorsque la construction cesse d'être évaluée comme une coordonnée
réelle et se réalise comme structure algébrique.
